\section{ОСНОВНАЯ ЧАСТЬ}

\subsection{Архитектура и взаимодействие компонентов}
Разработанная система управления счетами (\texttt{invoices-s13}) состоит из двух функциональных блоков, каждый из которых функционирует в изолированной среде:
\begin{enumerate}
    \item \textbf{Invoices Service} (Backend): Реализует основные операции над счетами (создание, чтение списка, поиск по ID). Взаимодействие осуществляется по протоколу gRPC на порту 8228. Сервис хранит данные в оперативной памяти (in-memory storage), что позволяет продемонстрировать логику работы без усложнения инфраструктуры внешними СУБД на данном этапе.
    \item \textbf{API Gateway} (Frontend Proxy): Выступает в роли посредника. Принимает внешние HTTP-запросы в формате REST (порт 8000), выполняет их трансляцию в gRPC-вызовы и возвращает результат клиенту в формате JSON.
\end{enumerate}

Схема информационных потоков:
\texttt{[External Client] -> (REST/JSON) -> [Gateway] -> (gRPC/Protobuf) -> [Invoices Service]}.

\subsection{Проектирование gRPC интерфейса}
Основой взаимодействия между сервисами является контракт, описанный на языке IDL (Interface Definition Language) в формате Protocol Buffers. Использование gRPC позволяет гарантировать целостность данных за счет строгой типизации полей.

\begin{lstlisting}[language=bash, caption=Контракт Invoices (invoices.proto)]
syntax = "proto3";
package invoices.v1;

service InvoicesService {
  rpc GetInvoice (GetInvoiceRequest) returns (GetInvoiceResponse);
  rpc CreateInvoice (CreateInvoiceRequest) returns (CreateInvoiceResponse);
  rpc ListInvoices (ListInvoicesRequest) returns (stream Invoice);
}

message Invoice {
  string id = 1;
  string name = 2;
  float amount = 3; // Индивидуальное поле варианта s13
}
\end{lstlisting}

\subsection{Реализация логики сервисов}
Сервис \texttt{invoices-svc-s13} реализован с использованием библиотеки \texttt{grpcio}. Особое внимание уделено реализации метода \texttt{ListInvoices}, использующего механизм серверного стриминга для эффективной передачи больших списков данных.

Шлюз реализован на базе асинхронного фреймворка FastAPI. Для каждой REST-операции создается gRPC-канал связи с бэкендом.

\begin{lstlisting}[language=Python, caption=Реализация метода создания счета в Gateway]
@app.post("/api/invoices", response_model=InvoiceSchema, status_code=201)
async def create_invoice(invoice: CreateInvoiceSchema):
    client = get_grpc_client()
    try:
        # Вызов gRPC метода бэкенда
        response = client.CreateInvoice(invoices_pb2.CreateInvoiceRequest(
            name=invoice.name, 
            amount=invoice.amount
        ))
        i = response.invoice
        return InvoiceSchema(id=i.id, name=i.name, amount=i.amount)
    except Exception as e:
        raise HTTPException(status_code=500, detail=str(e))
\end{lstlisting}

\subsection{Контейнеризация и инфраструктура}
Для каждого сервиса подготовлен \texttt{Dockerfile}, описывающий среду выполнения (Python 3.11-slim) и необходимые зависимости. Сборка и запуск всей системы автоматизированы через \texttt{docker-compose.yml}.

Использование Docker-сетей обеспечивает внутренний DNS-резолвинг: шлюз обращается к бэкенду по имени сервиса \texttt{invoices-svc-s13}, что соответствует принципам Service Discovery в распределенных системах.
